
\section{DESI survey}

In this section, we briefly present the Dark Energy Spectroscopic Instrument (DESI). This novel, fully robotic instrument, has allowed us to observe the Universe in great detail and 
proven to being capable of adding insight to our understanding of the accelerated Universe (\cite{DESI_DR2_2025}). 
In april 2026, DESI completed its main observation sequence and covered $14000\text{deg}^2$ with over 60 million objects (with 47 million galaxies and quasars) observed within 5 years. \\

\subsection{Spectroscopic instrument}

DESI is a fourth-generation spectroscopic survey. It is a robotic fibre-fed spectroscopic instrument. 
It is installed on the 4-meter Mayall telescope at Kitt Peak National Observatory in Arizona, USA. 
It contains 5000 fibres, each able to observe a separate spectrum, allowing the instrument to gather 5000 spectra at a time, covering for one exposure around $3 \text{deg}^2$
part of the sky. This is a great improvement over its predecessor SDSS (\cite{Dawson_2016}) which had 1000 fibres positioned manually, therefore taking a lot more time than DESI. 
The DESI spectrograph covers a range of $3600-9800 \AA$ where the light is divided into three wavelength channels: Blue ($b$), Red ($r$) and Near Infrared ($Z$). Finally, the light 
is captured by CCDs installed in a vacuum cryostat. Two different CCD models are used: STA4150 (4096 $\times$ 4096 pixels) and CCDs produced by LBL specifically for DESI 
(4114 $\times$ 4128 pixels). More details on the instrument can be found here \cite{DESI_Instrument_2016}.\\

\subsection{Tracers}

DESI catalogues its observations under multiple classes each presenting a of matter tracer in the observable Universe (\cite{DESI_Design_2016}).

\noindent \textbf{\underline{BGS}} \\

The first catalogue is the Bright Galaxy Survey (BGS). It is a flux limited galaxy catalogue that is divided into two parts based in the magntidue of the obsereved galaxy. BGS 
Bright contains galaxies with magnitudes in the $r$ band $<19.5$ and BGS faint with magnitudes $19.5 - 20.175$, these are present at a redshift range of $z<0.6$. To identify if the obsereved object is a galaxy or a star it is 
cross references with the Gaia catalogue. If it is not present then it is a galaxy. There are over 10 million galaxies in the BGS catalogue. These also present the galaxies closest to us. \\

\noindent \textbf{\underline{LRG}} \\

The second catalogue is the Luminous Red Galaxies (LRG) which contains old elliptical galaxies, identified due the $4000\AA$ break line in their spectra, in which 
the star formation rate has stopped. The LRG present galaxies in the redshift range $0.4<z<1.0$.

\noindent \textbf{\underline{ELG}} \\

The third catalogue is the Emission Line Galaxies (ELG) which contains younger spiral galaxies with star formation. It is divided into two subcategories: ELG VLOP with galaxies 
at redshifts $0.6<z<1.1$ and ELG LOP with galaxies at $1.1<z<1.6$. \\ 

\noindent \textbf{\underline{QSO}} \\

The fourth catalogue is the Quasar-Stellar Objects (QSO). The catalogue contains quasars, active galactic nuclei which are powered by the accretion of super massive black holes. 
The quasars measured at a redshift of $0.9<z<2.1$ are used as direct matter tracers, while quasars at $2.1<z<4.0$ will be used for their Lyman-$\alpha$ forest. This 
represents the furthest objects obsereved by DESI. \\

DESI at the time of writing this thesis has entered an extension phase of observation until 2028 where the sky coverage will go up to $17000\text{deg}^2$. \\

\subsection{ZTF and DESI Synergies}

ZTF and DESI present two of the current generation surveys that will tackle known tensions in the cosmological model. In addition to that, ZTF's sky coverage is in fact encompassed 
by DESI's coverage and the BGS sample is actually present at the same redshift range as ZTF. This creates synergies between the surveys that could explored.\\ 
Firstly, DESI's spectroscopic instrument has more accurate measurements of redshifts than the SEDM, therefore redshifts provided by DESI to ZTF can improve 
the mapping of the expanding history of the Universe using supernovae. They will also provide more accurate measurements of growth rate of structures 
using ZTF since it heavily relies on comparing the observed redshift to the cosmological redshift due to the expansion (see Section \ref{sec:vp_methods}). \\ 
Secondly, the BGS Bright sample contains galaxies in the same redshift range as the ZTF DR2. This highly favours conducting cross-survey analyses, which is the main 
goal of this PhD regarding the measurement of the growth rate of structures $f\sigma_8$. The DESI growth rate of structure measurements at low redshift is limited by the cosmic 
variance (the statistical uncertainty that arises due to only having access to one realization of the Universe) and galaxy bias. Supernova measurements can directly probe 
the underlying velocity field (see Section \ref{sec:vp_methods}) independently of any galaxy bias. Therefore a combination of both DESI and ZTF is projected to yield more 
precise measurements of $f\sigma_8$ at a low redshift. This is particularly important since the measurement of the growth rate of structures is a direct test of GR and is most 
sensitive at low redshift. \\
